\section{Discussion}
\label{sec:discussion}

\subsection{Findings}

\paragraph{F1: Completion, not prose quality, is the first-order obstacle.}
Four of ten models produced no learning graph, and four of the six that did did not finish stage~3. None of the stopped runs stopped because of how the prose read. The failures in Table~\ref{tab:failures} are failures to emit a well-formed tool call, to keep acting instead of announcing, to stay on task as the conversation fills the window, and to finish within a time limit. On this workload the ability that gates progress is protocol discipline over a long session, which is why the benchmark was framed around following a skill and calling tools and not around writing.

\paragraph{F2: A clean first stage did not predict a finished book.}
The cleanest learning graph came from the slowest model, which then wrote no chapter within the limit. The graph with the most duplicate labels, 40 of 200, came from the fastest model, which went on to write ten chapters. Stage~1 structural quality should not be used to choose a model for the whole pipeline, though with one run each this is an observation and not a rule.

\paragraph{F3: Sparse models were fastest, but speed did not buy completion.}
Both models that wrote every chapter activate about 3 billion parameters per token, and so does \Gtwentysix{} (3.8 billion), the fastest model, which wrote ten of fourteen. The largest sparse model, \Gptoss, ran at a usable speed and wrote six of ten. Architecture and speed are neither necessary nor sufficient for completion here; per-turn behavior, which the harness multiplies by hundreds of turns, matters more.

\paragraph{F4: Model-only tests did not predict harness behavior.}
\Gthirtyone{} passed every format rule in isolation and had the highest tool-call failure rate in the harness, 20\%. Model-only tests ran with reasoning off and the harness with each model's default, so the speeds measured in the first did not match the second. Two models wrote their first tool call as text in the harness, which may reflect a mismatch with the tool-call format the harness expects as much as an inability to call tools. The two layers of testing are not redundant; each finds failures the other cannot.

\paragraph{F5: Failures can be self-inflicted and can propagate across stages.}
Three of \Gtwentysix's chapter runs halted because the skill's own check found that a chapter depended on concepts placed in later chapters, and the chapter order that violated the dependency was the model's own stage~2 design. A cheap structural validation of the stage~2 output against the stage~1 graph would have found this in seconds, before ninety minutes of stage~3 per chapter. The benchmark chains stages by copying output forward, so a defect inherited from stage~2 is paid for once per chapter.

\paragraph{F6: The conformance score separates evenness from level, and shows its own limits.}
The headline separates two things that a mean would blur: \Gptoss{} is the most even book and the smallest, and \Nemo{} wrote the most and has the lowest floor. At the same time the score ties its top two books, cannot see that most of \Gptoss's words are specifications, and records that every book exceeds the target reading grade. It is a screen that points to chapters worth reading, not a verdict.

\paragraph{F7: On unified memory, context is a competitor to weights.}
\Mistral{} loaded and ran at 3 tokens per second and left no context of useful length. Two models drove a task to within about a thousand tokens of their configured window. One of the two units stopped responding twice during runs of the same model at large contexts. The cause was not established, and the second freeze at a smaller context makes memory exhaustion a weaker suspect than it first appeared, but the practical lesson holds: start a memory watchdog and a vitals log before any run, because a unified-memory machine can freeze instead of failing one process.

\subsection{Implications for Practitioners}

These follow from the observations above and are offered as practice, not as proof.

\begin{itemize}
\item \textbf{Gate on completion, then on quality.} Run a model-only test, then a three-minute harness smoke test, before any hour-long stage. A model that cannot call a tool in the harness will not finish a book.
\item \textbf{Budget time by turns, not tokens.} For the fast models the per-turn overhead exceeded generation time. Measure overhead per model; it ranged from 7 to 100 seconds here.
\item \textbf{Validate between stages.} Check each stage's output structurally before spending the next stage's budget on it.
\item \textbf{Treat a book-level score as a screen.} Read the chapters it flags, and do not choose a model on a headline that sits within the reweighting noise.
\item \textbf{Keep the evidence and the machine safe.} Isolate raw output, publish through a narrow path, and run the watchdog.
\end{itemize}

\subsection{Harness--Model Confound}

The harness is Claude Code, chosen because the published books were generated with it and the skills were written for it, so that the model is the only thing that changes. The cost is that its own instructions and tool definitions run to about 19,000 tokens and were written for Claude models. A local model that struggles may be struggling with the harness, not with the task. A second harness with a lighter prompt would separate a model failure from a harness mismatch, and is the most valuable single addition to this design.

\subsection{Threats to Validity}

\begin{itemize}
\item \textbf{One run per cell.} Run-to-run variance is unknown, and the ordering of models is not evidence. A second run could move a model from complete to incomplete or the reverse.
\item \textbf{No human judgment.} No person has judged whether any chapter is mathematically correct or pedagogically sound. The conformance score measures adherence to rules, and its thresholds were read off the same six books it now scores, so the score is not independent of the corpus.
\item \textbf{No cloud baseline.} The published reference book was made with earlier skills and no timings were kept, so the paper cannot say how a local model compares with a cloud model under identical conditions.
\item \textbf{Uncontrolled variables.} Context length, reasoning mode, and server behavior differed across models (Table~\ref{tab:controls}); one server's re-reading may have penalized one model for reasons that are not the model's.
\item \textbf{Time limits are choices.} A model that wrote nothing in two 90-minute runs might finish in 180 minutes; ``wrote no chapter'' is relative to the limit.
\item \textbf{A lenient definition of written.} A chapter counts as written with at least 1,000 words and no placeholder; the score then measures how well it conforms.
\item \textbf{Assistant-made decisions.} An AI assistant chose the candidate list, the model server, and the scoring methods under the author's direction. The design-choices page attributes each decision, but the assistant's blind spots are in the benchmark too.
\item \textbf{One machine, one course, one skill version.} Nothing here generalizes to another course, another skill version, or another hardware class without a rerun. Publisher claims about models were not checked.
\end{itemize}

\subsection{Responsible Use}

The opportunities and risks of large language models in education are well surveyed \cite{kasneci2023chatgpt}, and the central caution applies here with full force: generated educational content must be reviewed by a subject expert before it reaches learners. Local generation changes where data goes, not whether the text is correct. The book labels its generated content as unreviewed; its MicroSims are implemented but not approved; and its model-results pages are shown as each model wrote them. The conformance score should not be used to grade a person's work, and it was not designed for that.
